\Needspace{9\baselineskip}\pdfbookmark[1]{Göttingische gelehrte Anzeigen. 1811 Mai 25.}{prasse-1811}
\begin{center}
\rule{0.60\linewidth}{0.3pt}\\
Göttingische gelehrte Anzeigen. 1811 Mai 25.\\
\rule{0.60\linewidth}{0.3pt}
\end{center}

\textit{Logarithmische Tafeln für die Zahlen, Sinus und Tangenten, neu geordnet
von \textup{\textsc{Moritz von Prasse}}, ordentlichem Professor der Mathematik zu Leipzig.} 80 Seiten in Oktav. Leipzig. In Kommission bei P. J. Besson.

Diese Tafeln enthalten dasselbe wie die beliebten kleinen Tafeln von \textsc{Lalande}: die Logarithmen aller Zahlen bis 10000 und die Logarith\SourceLeaf{251}%
men der Sinus und Tangenten für jede einzelne Minute des Quadranten, alles auf fünf Dezimalstellen. Herr von \textsc{Prasse} hat dies jedoch bei einem nicht viel größeren Format auf ein Drittel der Seitenzahl reduziert. Dazu hat er die bei größeren Tafeln übliche Anordnung übernommen, jeweils zehn Logarithmen in eine Zeile zu setzen und die ersten Ziffern nur einmal anzugeben; dabei mussten allerdings sämtliche Differenzen wegfallen. So scheint an Bequemlichkeit wieder verloren zu gehen, was an Kürze gewonnen wird. Da sich dies jedoch nur nach tatsächlichem Gebrauch beurteilen lässt, hat der Rezensent, der an \textsc{Lalande}s kleine Tafeln gewöhnt war, diese eine Zeit lang beiseitegelegt und versucht, die vorliegenden zu benutzen. Er hat festgestellt, dass bei den Logarithmen der Zahlen, die hier 31, bei \textsc{Lalande} 111 Seiten einnehmen, der Vorteil, viel weniger blättern zu müssen, jene kleinen Unbequemlichkeiten deutlich überwiegt; daher benutzt er diese neuen Tafeln gern. Anders war seine Erfahrung bei den trigonometrischen Tafeln, die hier 40, bei \textsc{Lalande} 90 Seiten einnehmen. Dies liegt besonders daran, dass bei Herrn von \textsc{Prasse} die Bögen von 0 bis 90 Grad fortlaufen und dadurch die Sinus und Tangenten von den Kosinus und Kotangenten getrennt sind. Das ist umso lästiger, als man sehr häufig etwa den Sinus und Kosinus eines Bogens zugleich braucht oder, ohne den Bogen selbst zu benötigen, aus dem Sinus oder der Tangente den Kosinus bestimmen will. Hier würde er daher stets \textsc{Lalande}s Tafeln vorziehen. Er hätte gewünscht, dass Herr von \textsc{Prasse} stattdessen jede Seite in der Mitte noch einmal durch eine waagerechte Linie geteilt hätte. Damit wäre jene unangenehme Trennung vermieden worden, während die Unterbringung auf dem kleinen Raum dennoch möglich geblieben wäre.

Außerdem unterscheiden sich diese Tafeln dadurch, dass die letzte Ziffer jedes Logarithmus, wenn sie erhöht worden ist, in einer anderen Schrift gesetzt wird. Herr von P. glaubt, damit eine größere Rechengenauigkeit fördern zu können. Da man jedoch in der Praxis meist mit Logarithmen rechnen muss, die zu interpolieren sind, lässt sich jener Umstand nur umständlich berücksichtigen. Wenn man meint, dass Logarithmen, die nur bis auf eine halbe Einheit der fünften Dezimalstelle zuverlässig sind, keine hinreichend genauen Ergebnisse liefern können, tut man besser daran, größere Tafeln mit sechs oder sieben Dezimalstellen zu verwenden. Der Rezensent kann diese Anordnung daher nicht für sehr nützlich halten, obwohl sie, allgemein gesprochen, die Genauigkeit der \SourceLeaf{252}
Rechnung durchaus \emph{verdoppeln} kann, zumal die kursiven Ziffern neben den anderen dem Auge unangenehm und stellenweise nicht scharf genug sind. Sonst ist der Druck sauber; nur werden diejenigen, die solche Tafeln viel benutzen, stärkeres Papier wünschen.

